\documentclass[12pt,a4paper]{article}

\usepackage[margin=1in]{geometry}
\usepackage{graphicx}
\usepackage{amsmath}
\usepackage{booktabs}
\usepackage{hyperref}
\usepackage{titlesec}
\usepackage{setspace}

\titleformat{\section}{\large\bfseries}{\thesection}{1em}{}

\title{\textbf{Synora – AI-Driven EV Analytics}}
\author{
\textbf{Group Project Report} \\
Krishiv Gupta \\
Hardik Hathwal \\
Kabir Sharma \\
Syed Darain Qamar \\
\small Field: Generative AI / Machine Learning in Urban Planning
}
\date{}

\begin{document}

\maketitle
\onehalfspacing

\section{Problem Statement}

The rapid adoption of Electric Vehicles (EVs) has outpaced the strategic planning of charging station infrastructure, leading to two primary issues: \textbf{Charging Demand Congestion} in high-traffic administrative zones and \textbf{Infrastructure Underutilization} in others.

The goal of this project is to develop a predictive analytics system that forecasts \textbf{Occupancy (\%)} and \textbf{Charging Volume (kWh)} across urban charging networks. Accurate forecasting allows city planners to:

\begin{itemize}
    \item Optimize the placement of new charging piles.
    \item Manage grid load dynamically based on predicted demand.
    \item Reduce waiting times for EV users by identifying peak congestion windows.
\end{itemize}

\section{Data Description}

The project utilizes the \textbf{UrbanEV Dataset}, focusing on urban charging infrastructure in Shenzhen, China.

\textbf{Scope:}
\begin{itemize}
    \item 275 administrative zones (TAZs)
    \item 1,362 charging stations
    \item 17,532 charging piles
\end{itemize}

\textbf{Timeframe:} September 1, 2022 – February 28, 2023 (Hourly Resolution)

\textbf{Key Features:}
\begin{itemize}
    \item Usage Metrics: Occupancy (\%), Volume (kWh), Charging Duration
    \item Pricing: Electricity Price, Service Price
    \item Metadata: Zone Area, Piles Count, Latitude/Longitude
    \item Temporal Data: Hour, Weekday, Seasonal Indicators
\end{itemize}

\section{Exploratory Data Analysis (EDA)}

Key insights from the data:

\begin{itemize}
    \item \textbf{Temporal Stationarity:} Strong diurnal patterns with peaks during late-night and midday hours.
    \item \textbf{Data Integrity:} 0\% missing values across core time-series columns.
    \item \textbf{Congestion Variance:} Average occupancy ~17.87\%, while commercial zones exceed 80\%.
    \item \textbf{Correlation:} Strong relationship between Occupancy and Total Price, suggesting price elasticity.
\end{itemize}

\section{Methodology}

An ensemble learning framework was implemented to forecast dual targets: Occupancy and Volume.

\textbf{Models Used:}
\begin{itemize}
    \item Random Forest
    \item XGBoost
    \item LightGBM
\end{itemize}

\textbf{Feature Engineering:}
\begin{itemize}
    \item Lag Features: 1h, 3h, 6h, 12h, 24h, 168h
    \item Cyclical Encoding: Sine/Cosine transformations for Hour and DayOfWeek
    \item Moving Averages: 6h, 12h, 24h rolling means
\end{itemize}

\textbf{Training Strategy:}
Time-based split:
\begin{itemize}
    \item Training: September – January
    \item Testing: February
\end{itemize}

\section{Evaluation}

Models were evaluated using MAE, RMSE, and $R^2$ on the February 2023 test set.

\begin{table}[h]
\centering
\begin{tabular}{l l c c}
\toprule
\textbf{Target} & \textbf{Model} & \textbf{$R^2$ Score} & \textbf{MAE} \\
\midrule
Occupancy (\%) & LightGBM & 0.9999 & 0.0553 \\
Occupancy (\%) & Random Forest & 0.9999 & 0.0074 \\
Volume (kWh) & XGBoost & 0.9604 & 44.25 \\
Volume (kWh) & LightGBM & 0.9599 & 43.84 \\
\bottomrule
\end{tabular}
\end{table}

\textbf{Note:} The high $R^2$ for occupancy is influenced by the inclusion of 1-hour lag features.

\section{Optimization}

Performance improvements included:

\begin{itemize}
    \item \textbf{Hyperparameter Tuning:} Grid Search optimization (e.g., LightGBM at 800 estimators, 31 leaves).
    \item \textbf{Memory Optimization:} Dataset reshaped to long-format panel structure using \texttt{int8} and \texttt{float32}.
    \item \textbf{Cross-Variable Lags:} 24-hour volume lag added to improve occupancy prediction.
\end{itemize}

\section{Team Contributions}

\begin{itemize}
    \item \textbf{Krishiv Gupta (Project Lead):} Identified and sourced the UrbanEV dataset; implemented and trained machine learning models.
    \item \textbf{Hardik Hathwal:} Led data preprocessing, cleaning, and transformation pipeline; assisted in evaluation metric computation.
    \item \textbf{Kabir Sharma:} Contributed to data preprocessing, feature preparation, and performance evaluation.
    \item \textbf{Syed Darain Qamar:} Designed and developed the user interface (UI) for analytics visualization; structured and authored the final project report.
\end{itemize}

\section{Conclusion}

Synora demonstrates how AI-driven predictive modeling can enhance urban EV infrastructure planning. By leveraging ensemble learning, feature engineering, and temporal modeling, the system enables intelligent decision-making for congestion management, grid optimization, and infrastructure expansion.

\end{document}